\documentclass[10pt,a4paper]{amsart}
\usepackage[margin=19mm]{geometry}
\usepackage{amsmath,amssymb,mathtools}
\usepackage[scr=rsfs,cal=euler]{mathalfa}
\usepackage{xfrac}
\usepackage[unicode,hidelinks,bookmarks=false]{hyperref}
\newcommand{\ie}{i.e.\ }
\newcommand{\cf}{cf.\ }
\newcommand{\resp}{respectively\ }
\newcommand{\Subsection}{\subsection}
\providecommand{\nobreakdash}{\nobreak\hbox{-}}
\setlength{\emergencystretch}{2em}
\multlinegap0pt
\usepackage{xcolor}
\newtheorem{theorem}{Theorem}
\newtheorem{corollary}{Corollary}
\newtheorem{proposition}{Proposition}
\newcommand{\R}{{\mathbf R}}
\newcommand{\re}{\color{red}}
\title{An elliptic proof of the splitting theorems from Lorentzian geometry}
\author{Mathias Braun, Nicola Gigli, Robert J. McCann, Argam Ohanyan, Clemens S\"amann}
\date{Source: arXiv:2410.12632v2 (CC BY 4.0)}
\begin{document}
\maketitle

\noindent\emph{Source and licence.} An elliptic proof of the splitting theorems from Lorentzian geometry. Version arXiv:2410.12632v2 (CC BY 4.0), distributed by arXiv under the Creative Commons Attribution 4.0 International licence, http://creativecommons.org/licenses/by/4.0/, as recorded in the arXiv licence metadata for this exact version and shown on the abstract page https://arxiv.org/abs/2410.12632v2. Full licence notice: \texttt{licenses/arxiv-2410.12632-CC-BY-4.0.txt}.\par\smallskip

\noindent\textit{Editorial scope.} Counted unit: the article's corollary C:asymptotes (source0.tex lines 1284--1289). The article's complete original proof of it is delivered with it (source0.tex lines 1291--1295, one \texttt{proof} environment of 5 lines). No mathematics is written, repaired or re-derived: the statement and the proof are the article's own lines, in the article's own notation, and no retained line is substituted or re-flowed.  Retained uncounted as the source-local context the proof needs: lines 196--205; lines 304--306; lines 675--683; lines 1242--1256.  Nothing was omitted from any retained span, so the package carries no omission record.  No retained line contains a citation, so the wrapper carries no bibliography and no bibliography entry.  No retained line contains a figure environment, an image-inclusion command or drawing code, so no image is delivered and the package declares no asset dependency.  Editorial formatting only, in addition: an AMS \texttt{amsart} preamble that restates the article's own declarations and macros used by the retained lines, namely the environments \texttt{theorem}, \texttt{corollary}, \texttt{proposition} on separate counters, together with the standard packages those lines need.  The retained environments print in this document as Theorem 1, Corollary 1, Proposition 1 in order of appearance, whereas the article prints the same items with the numbering of its own full document.  The complete native source of the article as distributed in the arXiv e-print for this exact version is bundled unchanged as \texttt{sources/166941/source0.tex}.
\begin{theorem}[Lorentzian splitting theorem]
\label{T:Lorentzian splitting}
If a spacetime $(M,g)$ satisfies the strong energy condition 
\begin{equation}\label{strong energy}
R(v,v) \ge 0\ \mbox{\rm for all timelike vectors}\ v,
\end{equation} 
contains a proper-time maximizing, isometrically embedded copy $\gamma$ of the Euclidean line $\R$,
and is either   (a) globally hyperbolic or (b) timelike geodesically complete,
then $(M,g)$ is isometric to $(\R \times S, dt^2 - h)$ where $(S,h)$ is a complete Riemannian manifold with nonnegative Ricci curvature.
\end{theorem}

\begin{equation}\label{b^+_r}
b^+_r(x) = \ell(\gamma(0),\gamma(r))-\ell(x,\gamma(r)); %=r-\ell(x,\gamma(r));
\end{equation}

\begin{corollary}[Busemann function $b^+$ is $p$-superharmonic]
\label{L:harmonicity}
Fix $0 \ne p<1$. 
Under the hypotheses of Theorem \ref{T:Lorentzian splitting}(a) or (b) on $(M,g)$ and $\gamma$,
if $V \subset M$ is a domain on which the functions $\{b^+_r\}_{r \ge R}$ of \eqref{b^+_r} are equi-semiconcave when $R$ is sufficiently large, and if in addition %{\re (DELETE: in case (b))} 
for all $p \in V$ and $r \ge R$ there exists a timelike maximizing geodesic from $p$ to $\gamma(r)$,
then $b^+ = \lim_{r\to \infty} b^+_r$ 
is semiconcave, $p$-superharmonic and $|db^+|=1$ on $V$.
\end{corollary}

\begin{proposition}[Equi-semiconcave Busemann limits on asymptotes]\label{P:asymptotes}
Let $(M,g)$ be an (a) globally hyperbolic 
spacetime.
Let $b^+=\lim_{r \to \infty} b^+_r$ denote the Busemann function associated by \eqref{b^+_r} to a future-directed ray $\gamma:[0,\infty)\longrightarrow M$.
Let $\alpha:[0,a_+)\longrightarrow M$ be a 
(future inextendible)
timelike asymptote to $\gamma$. %{\re 
Let $b^+ \in C^1(W)$ %{\re and $|db^+|=1$?} 
on a neighbourhood of $W$ of %$\beta$. If
$x_0=\alpha(a)$ for some $a \in [0, a_+)$. 
Then for large enough 
$C=C(M,g,\tilde g, x_0,db^+(x_0))$ and 
$R$ depending on the same parameters as well as on $\gamma$,
 the functions $(b^+_r)_{r \ge R}$ have semiconcavity constant $C$ on some smaller neighbourhood $X \subset I^-(\gamma(R))$ of $x_0$.
\end{proposition}

\begin{corollary}[$p$-superharmonicity of $b^+$ near asymptotes]
\label{C:asymptotes}
Assume the (a) globally hyperbolic spacetime $(M,g)$ satisfies the strong energy condition \eqref{strong energy}. Then 
inside the neighbourhood $X$ of $\alpha(a)$ identified in Proposition~\ref{P:asymptotes},
$b^+$ is $p$-superharmonic, semiconcave, and $|db^+|=1$.
\end{corollary}

\begin{proof}
Proposition \ref{P:asymptotes} provides the equi-semiconcavity of $\{b^+_r\}_{r \ge R}$ 
necessary for Corollary~\ref{L:harmonicity} to imply $|db^+|=1$ and semiconcavity and $p$-superharmonicity of $b^+$  on the interior of $X$;
global hyperbolicity ensures timelike maximizing geodesics connect each $x \in X \subset I^-(\gamma(R))$ to $\gamma(r)$ for $r \ge R$. 
\end{proof}

\end{document}
